\chapter{Bunions}

\section*{Definition and Overview}
A bunion is a bony bump that develops on the inner side of the foot at the base of the big toe. It forms when the great toe angles towards the second toe, a deformity known as hallux valgus, and the joint at its base (the first metatarsophalangeal joint) becomes displaced and prominent. A bunion is not a growth or tumour; it is a structural misalignment of the bones of the forefoot that usually develops gradually over many years.

Bunions can cause pain, redness, and difficulty fitting into shoes, and in some people they lead to arthritis of the joint. They are common and treatable, and many people manage their symptoms successfully without surgery. Treatment is tailored to the severity of the deformity, the level of symptoms, and the person's goals and activities.

\section*{Epidemiology}
Bunions are very common, particularly in women, who are affected several times more often than men. The condition becomes more frequent with advancing age and is uncommon before the teenage years. Studies suggest that roughly one in three adults has some degree of hallux valgus, although many remain free of symptoms. Rates are higher among people who habitually wear narrow, pointed, or high-heeled footwear, and the deformity shows a strong family tendency, often running in several generations of a family. The wide global variation in prevalence reflects the strong influence of footwear habits and other environmental factors.

\section*{Aetiology and Pathophysiology}
The exact cause of bunions is not fully understood, but a combination of inherited and environmental factors contributes:

\begin{itemize}
    \item \textbf{Footwear:} narrow, tight, or high-heeled shoes squeeze the toes together and are strongly associated with bunions, particularly in women
    \item \textbf{Heredity:} foot shape, joint flexibility, and flat-arched feet run in families and influence risk
    \item \textbf{Foot structure:} excessive pronation of the foot, hypermobile joints, and abnormal pull of the tendons around the great toe
    \item \textbf{Arthritis:} inflammatory joint disease such as rheumatoid arthritis can produce or accelerate the deformity
    \item \textbf{Injury:} trauma to the foot occasionally precipitates bunion formation
\end{itemize}

As the big toe is pushed towards the second toe, the balance of tendons and ligaments around the joint is disturbed. The first metatarsal drifts medially, causing the joint to protrude as a bony bump, while the proximal phalanx deviates laterally. Over time the joint surfaces become progressively misaligned, the protective cartilage may wear down, and secondary osteoarthritis can develop, contributing to pain and stiffness.

\section*{Clinical Features}
\begin{itemize}
    \item \textbf{Bony bump:} a visible prominence on the inner side of the foot at the base of the great toe
    \item \textbf{Pain and tenderness:} aching over the bunion, worsened by standing, walking, or tight footwear
    \item \textbf{Redness and swelling:} local inflammation, often worse after activity or prolonged shoe wear
    \item \textbf{Toe deviation:} the great toe leaning towards, overlapping, or being pushed beneath the second toe
    \item \textbf{Footwear difficulty:} rubbing against shoes, with calluses or corns over the toes
    \item \textbf{Sensory changes:} burning, numbness, or tingling over the joint in some people
    \item \textbf{Stiffness:} reduced movement of the big toe in advanced cases
\end{itemize}

\section*{Differential Diagnosis}
\begin{itemize}
    \item \textbf{Gout:} sudden, intensely painful, red, and swollen joint at the base of the big toe, often with elevated uric acid
    \item \textbf{Hallux rigidus:} osteoarthritis of the first metatarsophalangeal joint causing pain and stiffness without the typical outward deviation
    \item \textbf{Septic arthritis:} infection of the joint, usually with severe pain, fever, and limited movement
    \item \textbf{Rheumatoid arthritis:} symmetrical inflammatory joint involvement that may affect the forefoot
    \item \textbf{Ganglion cyst:} a soft-tissue swelling over the joint rather than a true bony deformity
    \item \textbf{Bursitis:} isolated inflammation of the bursa over the joint, without significant bone malalignment
    \item \textbf{Metatarsalgia or stress fracture:} pain under the forefoot from overload or fracture rather than joint deformity
\end{itemize}

\section*{When to Seek Medical Care}
People should see a doctor or podiatrist if a bunion causes persistent pain, interferes with daily activities or exercise, becomes increasingly red, hot, or tender, or is accompanied by new numbness or difficulty moving the toe. Earlier review is also worthwhile for people with diabetes or poor circulation, because foot problems can progress quickly in these groups.

\textbf{Call 000 (or local emergency services) for:}
\begin{itemize}
    \item Sudden, severe pain and swelling at the base of the big toe with fever or chills, which may indicate a septic joint
    \item Signs of spreading infection, such as red streaks ascending the foot or leg
    \item Sudden inability to move the toe or bear weight after an injury
\end{itemize}

\section*{Diagnosis and Investigations}
\begin{itemize}
    \item \textbf{History:} duration of the bunion, character of the pain, footwear habits, family history, occupation, and any other joint symptoms are explored
    \item \textbf{Examination:} inspection of both feet, assessment of toe alignment, range of joint movement, overlying skin changes, and circulation and sensation
    \item \textbf{X-ray:} weight-bearing films of the foot are the principal investigation, used to measure the hallux valgus and intermetatarsal angles and to assess joint degeneration
    \item \textbf{Blood tests:} inflammatory markers, uric acid, or autoimmune serology are requested only when gout or inflammatory arthritis is suspected
    \item \textbf{Ultrasound or MRI:} occasionally used to assess soft tissues, bursae, or suspected stress fracture
\end{itemize}

\section*{Management}
Most bunions are managed without surgery:
\begin{itemize}
    \item \textbf{Footwear:} shoes with a wide, deep toe box and a low heel are the single most helpful measure
    \item \textbf{Orthotics:} shoe inserts, toe spacers, and night splints may reduce pressure and slow progression
    \item \textbf{Pain relief:} simple analgesics such as paracetamol or non-steroidal anti-inflammatory medicines during flares
    \item \textbf{Ice and padding:} ice packs, bunion pads, and silicone sleeves cushion the prominence
    \item \textbf{Foot exercises:} toe stretches and strengthening exercises help maintain joint mobility
    \item \textbf{Podiatry:} professional advice, padding, and custom orthoses from a podiatrist
    \item \textbf{Surgery:} bunionectomy is considered when pain persists despite conservative care and significantly impairs daily life
\end{itemize}

\section*{Complications}
\begin{itemize}
    \item Bursitis over the joint from friction against footwear
    \item Calluses, corns, and skin irritation of the adjacent toes
    \item Osteoarthritis of the first metatarsophalangeal joint over time
    \item Secondary deformity of the lesser toes, including hammering or dislocation
    \item Pain that limits walking, exercising, or working
    \item Complications of surgery if undertaken, including recurrence, infection, nerve irritation, or stiffness
\end{itemize}

\section*{Prognosis and Outcomes}
Bunions tend to progress slowly over many years, but a substantial proportion of people remain comfortable with footwear adjustments and activity modification, since symptoms may come and go. Where pain is persistent and conservative measures have failed, surgery provides good relief in the majority of cases, although the deformity can recur in a minority of people. Attention to footwear early in the course may slow progression, especially in younger people and those with a strong family history.

\section*{Prevention and Self-Care}
\begin{itemize}
    \item Choose shoes with a wide, deep toe box and a low heel, and buy shoes late in the day when the feet are largest
    \item Avoid narrow, pointed, or high-heeled shoes for everyday wear
    \item Wear well-fitting, cushioned socks and inspect the feet regularly for rubbing or redness
    \item Use bunion pads or toe spacers to reduce pressure
    \item Keep the foot flexible with simple toe stretches and strengthening exercises
    \item Maintain a healthy weight to reduce load through the forefoot
    \item Seek podiatry advice early if problems develop or progress
\end{itemize}

\section*{Sources and Further Reading}
The following sources provide reliable, accessible information on bunions and hallux valgus:
\begin{itemize}
    \item NHS. \textit{Bunions.} https://www.nhs.uk/conditions/bunions/
    \item Mayo Clinic. \textit{Bunions.} https://www.mayoclinic.org/
    \item MedlinePlus. \textit{Bunions.} https://medlineplus.gov/
    \item MSD Manual. \textit{Hallux valgus and hallux rigidus.} https://www.msdmanuals.com/
    \item American College of Foot and Ankle Surgeons. \textit{Bunions.} https://www.acfas.org/
    \item NICE. \textit{Hallux valgus clinical guidance.} https://www.nice.org.uk/
\end{itemize}
